\documentclass[11pt]{article}
\usepackage[a4paper,margin=2.2cm]{geometry}
\usepackage[section]{placeins}
\usepackage[hidelinks]{hyperref}
\usepackage{stickyflow}

\title{\bfseries Figures that answer back}
\author{A sample paper written with the stickyflow package for LuaLaTeX}
\date{}

\begin{document}
\maketitle

\begin{abstract}
Every figure in this paper is written in SFL, a language that describes a visual artifact in text: its parts, what they mean, how
they relate, where they sit and how they look. The figures are drawn while the document compiles, in the fonts of the document, and the
PDF carries their source. In print they are ordinary figures. Opened in a viewer that reads SFL, they come alive: a slider changes a
distribution, a section is turned, an instrument runs. The paper shows charts, a field of data in three dimensions, instruments,
drawings that compute their own measures, and an object in three dimensions.
\end{abstract}

\section{Introduction}

A figure in a paper is usually the end of a process that happened somewhere else. The data was plotted by one program, the drawing
made in another, and what arrives in the document is a picture. Whoever reads the paper cannot ask the figure anything: not what lies
under a point, not what happens if a parameter changes.

The figures here take another way. Each one is a short text in SFL, kept with the document. The renderer that draws it in a browser
also draws it in the PDF, so both look alike. The text travels inside the PDF, and a viewer that knows SFL draws the figure again and
lets it react. The printed page stays an ordinary page.

\section{Data at a glance}

Figure~\ref{fig:charts} shows six common forms side by side: a series over time, a ranking, a cloud of points, the spread of a
measure, a matrix of scores and the density of two measures. Each is a few dozen lines of SFL with its data written in.

\begin{figure}[htbp]
\centering
\sflinput[0.32\textwidth]{line.sfl}\hfill
\sflinput[0.32\textwidth]{bars.sfl}\hfill
\sflinput[0.32\textwidth]{scatter.sfl}\\[6pt]
\sflinput[0.32\textwidth]{boxplot.sfl}\hfill
\sflinput[0.32\textwidth]{heatmap.sfl}\hfill
\sflinput[0.32\textwidth]{hexbin.sfl}
\caption{Six forms of data. Top: a series over a year against the year before, a ranking by month, a cloud of points colored by
its position. Bottom: visitors per day as boxes with every point, focus scores of six teams over eight weeks, and the density of two
measures in hexagonal bins.}
\label{fig:charts}
\end{figure}

Flows and parts of a whole are read better in other forms. Figure~\ref{fig:flow} follows energy from its sources through the grid to
where it is used. Figure~\ref{fig:parts} compares profiles and shares.

\begin{figure}[htbp]
\centering
\sflinput[0.75\textwidth]{sankey.sfl}
\caption{Where the power comes from and where it goes. The width of every band is the energy it carries.}
\label{fig:flow}
\end{figure}

\begin{figure}[htbp]
\centering
\sflinput[0.32\textwidth]{radar.sfl}\hfill
\sflinput[0.32\textwidth]{polar.sfl}\hfill
\sflinput[0.32\textwidth]{rings.sfl}\\[6pt]
\sflinput[0.49\textwidth]{treemap.sfl}\hfill
\sflinput[0.49\textwidth]{mekko.sfl}
\caption{Profiles and shares. Top: a team profile today and as targeted, strengths as a polar area, and three goals of a day as
rings. Bottom: a month of spending as nested areas, and market shares by category, where the width of a column is the size of its
category.}
\label{fig:parts}
\end{figure}

\section{A field in three dimensions}

Some data has a place. Figure~\ref{fig:field} shows the density of a mixture of three normal distributions on a grid of 36 by 36
cells, as a field of columns whose height and color are the density. In print the field is a still picture. In the viewer it can be
turned, and two sliders change the spread of the three populations and the height of the columns.

\begin{figure}[htbp]
\centering
\sflinput[0.8\textwidth]{density-skyline.sfl}
\caption{Three populations, one field. The camera takes in a room as high as the highest column can grow, so the view keeps its
size while the columns rise and fall.}
\label{fig:field}
\end{figure}

\section{Instruments}

An interface is a figure too. The three watch faces in Figure~\ref{fig:watches} are written in SFL, with their data sources and the
shaders that fill them. Shaders only run in a browser, so print shows a picture of each face, and the viewer shows the face itself,
running.

\begin{figure}[htbp]
\centering
\sflinput[0.27\textwidth][][weather-watch.png]{weather-watch.sfl}\hfill
\sflinput[0.27\textwidth][][plasma-watch.png]{plasma-watch.sfl}\hfill
\sflinput[0.27\textwidth][][halftone-watch.png]{halftone-watch.sfl}
\caption{Three watch faces: the weather over a moving sea, a battery as a plasma gauge, and the time in halftone dots.}
\label{fig:watches}
\end{figure}

\section{Drawings that compute}

A construction drawing is a set of measures. When the measures are written as expressions, the drawing follows them.
Figure~\ref{fig:letter} builds a capital letter on a grid of 50 units, with its metrics, its circles, the handles of its curves and four
steps from the stem to the leg.

\begin{figure}[htbp]
\centering
\sflinput[0.62\textwidth]{letter.sfl}
\caption{How a capital is built: cap height, bowl and baseline, the circles of the bowl and the counter, and the angle of the leg.}
\label{fig:letter}
\end{figure}

Figure~\ref{fig:house} is a study of a house on rails: three buildings in a row and a shell of larch that rides over them. Plan,
elevation and section come from one set of parameters. In the viewer the shell slides along its rails, the doors of the glasshouse
open, the car leaves the garage, and the hour moves the sun, the shadows and the light in the windows.

\begin{figure}[htbp]
\centering
\sflinput[\textwidth]{sliding-house.sfl}
\caption{A study after a house with a sliding shell, drawn from its parameters: ground floor plan, north elevation on the axes of the
plan, and a section across.}
\label{fig:house}
\end{figure}

\section{An object in three dimensions}

The last figure is an instrument of another kind: a synthesizer, built from solids with keys that go down and a panel drawn on its
case. Its models are attached to the PDF, so the viewer can draw it and let it be turned.

\begin{figure}[htbp]
\centering
\sflinput[0.9\textwidth][][juno.png]{juno.sfl}
\caption{A synthesizer of the eighties in three dimensions. In print a still picture of the scene; in the viewer the scene itself.}
\label{fig:juno}
\end{figure}

\section{Conclusion}

Written as text, a figure can be read, reviewed and changed like the rest of a paper, and it keeps the meaning of its parts. Drawn by
the same renderer in print and on screen, it looks the same in both. Carried inside the PDF, it can be asked questions long after the
paper was printed.

\end{document}
